% Competition-paper LaTeX source driven by formal repository results.
% Build with: python scripts/build_paper.py
\documentclass[UTF8,12pt,a4paper]{ctexart}

\usepackage[a4paper,margin=2.2cm]{geometry}
\usepackage{amsmath,amssymb,bm}
\usepackage{booktabs,tabularx,longtable,multirow}
\usepackage{graphicx,float,placeins}
\usepackage{hyperref}
\usepackage{xcolor}
\usepackage{enumitem}
\usepackage{listings}
\usepackage{caption}

\hypersetup{colorlinks=true,linkcolor=black,citecolor=black,urlcolor=blue}
\setlength{\parindent}{2em}
\setlength{\parskip}{0.15em}
\renewcommand{\topfraction}{0.90}
\renewcommand{\bottomfraction}{0.85}
\renewcommand{\textfraction}{0.08}
\renewcommand{\floatpagefraction}{0.70}

\lstset{
  language=Python,
  basicstyle=\small\ttfamily,
  breaklines=true,
  frame=single,
  numbers=left,
  numberstyle=\tiny
}

% Generated by scripts/generate_latex_assets.py
\input{generated/result_macros.tex}

\title{多源融合机器人定位及任务优化}
\author{}
\date{}

\begin{document}
\maketitle

\begin{abstract}
本文面向异频、异步且含噪的双源机器人定位数据，建立从时间同步、相对系统偏差辨识、异步状态融合到固定轨迹任务调度的统一建模链。为避免各问内部代码符号约定不一致，全文统一采用
\[
  t_{2,\mathrm{aligned}}=t_2+\delta
\]
报告时间偏差。问题一得到 $\delta=\QOneDelta\,\mathrm{s}$，对齐 RMSE 为 $\QOneRMSE\,\mathrm{m}$；问题二得到 $\delta=\QTwoDelta\,\mathrm{s}$ 以及相对空间偏差 $(\QTwoBiasX,\QTwoBiasY)\,\mathrm{m}$，并采用异步 Kalman 滤波与 RTS 平滑恢复 10 Hz 轨迹；问题三通过 HAC--Wald、移动块 bootstrap 与工程效应阈值联合判定未发现需要显式建模的系统偏差，统一时间偏差为 $\delta=\QThreeDelta\,\mathrm{s}$；问题四在固定融合轨迹上构造连续可行时间窗和候选任务，使用三阶段字典序 MILP 在模板九行容量内完成 \QFourTaskCount 项任务，其中射击 \QFourShootCount 项、拍照 \QFourPhotoCount 项。MILP 最小归一化安全裕度为 \QFourMinMargin，约为贪心方案的 \QFourMarginRatio 倍。所有正式数值、表格和图件均由仓库结果文件自动生成并通过独立审计脚本检查。

\noindent\textbf{关键词：}多源定位；时间对齐；Kalman 滤波；RTS 平滑；HAC--Wald；混合整数线性规划
\end{abstract}

\tableofcontents
\clearpage

\section{问题重述与统一符号}

机器人由两类定位系统给出异频观测。前三问依次要求在无噪声、含噪且含固定系统偏差、实际复杂误差三种条件下恢复统一 10 Hz 轨迹；第四问在固定轨迹上安排射击和拍照任务。

为了避免不同代码模块对“时间偏差”的正负号定义不一致，本文只在报告层使用一个符号：
\begin{equation}
  t_{2,\mathrm{aligned}}=t_2+\delta.
\end{equation}
因此 $\delta>0$ 表示方式 2 时间戳整体增加，$\delta<0$ 表示整体减小。底层程序仍保留各自已经验证的内部变量定义，自动报告脚本负责转换到统一口径。

\input{generated/table_time_offsets.tex}

四问的总体技术路线为
\[
\text{时间同步}
\longrightarrow
\text{时空标定}
\longrightarrow
\text{偏差检验与鲁棒融合}
\longrightarrow
\text{固定轨迹任务调度}.
\]

所有正式结果均写入 \texttt{05\_results/}，论文图来自 \texttt{06\_figures/}，因此论文正文不再作为数值的唯一来源。

\section{问题一：无噪声异步数据的时间对齐}

设两类定位数据均为同一条连续真实轨迹的无噪声离散采样。先分别构造连续轨迹插值，在公共时间域上以运动特征互相关给出多个粗候选，再用位置域均方误差
\begin{equation}
J(\delta)=\frac{1}{N(\delta)}\sum_{t_k\in\Omega(\delta)}
\left\|\bm r_1(t_k)-\bm r_2(t_k-\delta)\right\|_2^2
\end{equation}
进行全局粗扫与有界连续精修。互相关只用于缩小候选范围，最终时间偏差由位置 MSE 决定，从而降低重复轨迹形状产生伪相关峰的风险。

正式结果为 $\delta=\QOneDelta\,\mathrm{s}$，对齐 RMSE 为 $\QOneRMSE\,\mathrm{m}$，达到数值误差量级。对齐后将两类无噪声采样点合并，并在不跨异常大时间空洞的前提下重建 10 Hz 轨迹。

问题一额外通过多插值器比较、已知时间偏差合成试验、搜索范围与步长敏感性分析验证结果稳定性。由于该问题不存在随机测量噪声，近零对齐残差是模型正确性的直接检验之一。

\section{问题二：稳健时空标定与异步状态融合}

问题二同时存在随机测量噪声、时间不同步和固定相对空间偏差。设两源观测为
\begin{align}
\bm z_1(t_i)&=\bm p(t_i)+\bm\varepsilon_{1,i},\\
\bm z_2(s_j)&=\bm p(s_j+\delta)+\bm b+\bm\varepsilon_{2,j},
\end{align}
其中 $\bm b$ 只表示两套设备之间的相对偏差。在无外部绝对真值时，不把它归因于任一单独设备。

模型先用偏差不敏感运动特征做粗配准，再在固定公共评价网格上建立 Huber 剖面目标，联合估计时间偏差和二维常偏差。固定评价样本集可避免候选时间偏差改变公共样本集合而人为移动目标函数。

校正后不把 10 Hz 插值点作为新观测，而是在原始 4 Hz/5 Hz 异步事件时刻运行常加速度 Kalman 滤波，并用 99\% NIS 门控对异常观测稳健降权，再采用 RTS 固定区间平滑恢复完整状态。过程噪声在关闭门控反馈的调参阶段根据门控前 NIS 和 Cholesky 白化创新相关性选择。

正式估计为统一时间偏差 $\delta=\QTwoDelta\,\mathrm{s}$，相对系统偏差约为 $(\QTwoBiasX,\QTwoBiasY)\,\mathrm{m}$，偏差模为 \QTwoBiasNorm\,m。校正后两源位置差 RMSE 为 \QTwoCorrectedRMSE\,m。

\input{generated/table_q2.tex}

\section{问题三：偏差检验驱动的稳健融合}

问题三首先估计时间偏差，再判断两类实际观测之间是否存在需要显式建模的系统偏差。为处理时间相关残差，本文不采用独立同分布假设下的普通均值检验，而使用 HAC 协方差构造二维 Wald 统计量，并辅以移动块 bootstrap 置信区间和 HAC 趋势检验。

设对齐后的两源位置差为 $\bm d_t$，检验
\[
H_0:\ E(\bm d_t)=\bm 0.
\]
在显著性判断之外，引入偏差模相对于稳健噪声尺度的工程效应指数。只有当统计显著且工程效应达到预设阈值时，才启用系统偏差状态；若同时检出显著趋势，再允许偏差随机游走。

正式数据下统一时间偏差为 $\delta=\QThreeDelta\,\mathrm{s}$，HAC--Wald $p=\QThreeWaldP$，偏差模约 \QThreeBiasNorm\,m，工程效应指数为 \QThreeEffectIndex。移动块 bootstrap 的两个坐标区间均覆盖 0，趋势检验亦未发现显著漂移，因此正式融合不启用系统偏差状态。

随后在校正后的原始异步观测事件上运行稳健 Kalman 滤波和 RTS 平滑，并传播状态协方差得到 10 Hz 位置、速度、加速度及模型内位置标准差。该结论应表述为“当前样本与阈值下未发现具有统计和工程意义的相对系统偏差”，而不是证明真实偏差严格等于零。

\input{generated/table_q3.tex}

\section{问题四：固定轨迹上的时空任务调度}

问题四不再优化机器人路径，而是在问题三的固定融合轨迹上选择“何时、对哪个目标、执行什么任务”。由平滑轨迹计算位置、速度、加速度以及机器人到目标的距离和方位角，并在 0.1 s 候选网格上检查完整准备窗口。

射击候选要求准备窗口内持续满足
\[
5\le d\le30,\qquad v\le2,\qquad a\le1.5,
\]
且准备时间为 1.5 s；拍照候选要求
\[
10\le d\le40,\qquad v\le1.5,\qquad a\le1.5,
\]
并按 0.5 s 对准窗口处理。同一拍照目标的任意两次有效照片满足圆周最小方位角差不小于 $60^\circ$。候选生成后再用 0.01 s 连续状态评估进行复核，避免 10 Hz 离散网格遗漏窗口内部的约束越界。

最终调度采用三阶段字典序 MILP：第一阶段在结果模板九行容量内最大化任务数；第二阶段固定任务数后最大化所有已选任务的最小归一化硬约束裕度；第三阶段保持前两级最优值，进一步最大化总裕度与拍照角度多样性。正式方案完成 \QFourTaskCount 项任务，其中射击 \QFourShootCount 项、拍照 \QFourPhotoCount 项，射击期望命中数为 \QFourExpectedHits。MILP 最小裕度为 \QFourMinMargin，而最早可行贪心基线仅为 \QFourGreedyMargin，前者约为后者的 \QFourMarginRatio 倍。

题面未明确说明不同任务的准备过程是否允许并行。正式主模型采用“单机器人同一时段只准备/执行一项任务”的保守工程假设，但该假设不作为题面原文引用。仓库同时运行允许准备阶段并行的对照模型，并将两种资源模式的一级最优任务数与安全裕度写入独立结果文件，从而判断该附加假设是否改变核心结论。

\input{generated/table_q4_summary.tex}
\input{generated/table_q4_schedule.tex}

\section{模型验证、稳健性与工程交付}

本项目将“数值求解正确”与“提交材料正确”分开审计。数值层面，Q1 通过精确公共坐标匹配和合成偏差恢复验证时间同步；Q2/Q3 通过 NIS、白化创新相关、协方差正定性、敏感性分析和合成数据检验状态融合；Q4 对所有最终两位小数任务时刻重新以 0.01 s 细网格检查完整准备区间，并对模板未写区域做值与样式快照比较。

交付层面新增三类自动检查：
\begin{enumerate}[label=(\arabic*)]
  \item \textbf{输入审计：}将本地附件的 SHA256、字节数、sheet 名、行数与字段和 \texttt{00\_problem/input\_manifest.json} 对照；
  \item \textbf{结果审计：}检查 Q1--Q3 轨迹的有限性、严格递增和 10 Hz 步长，检查统一时间偏差转换、Q4 模板容量和 MILP gap，并确认每个正式图件清单中的 PDF 存在；
  \item \textbf{提交审计：}在打包前生成 JSON/Markdown 报告和 SHA256 文件清单，避免把历史论文、临时文件或错误版本的 \texttt{result.xlsx} 混入最终 ZIP。
\end{enumerate}

论文数值通过 Python 自动生成的 TeX 宏与表格注入，而不是在多个章节手工重复录入。因此当正式结果文件变化时，重新运行资产生成器即可同步论文，降低“代码结果已更新、正文仍保留旧数字”的风险。

需要强调的是，模型内位置标准差和 Q4 扰动试验反映的是给定状态空间模型及压力情景下的稳定性，不应被解释为覆盖全部模型不确定性的严格概率保证。系统偏差结论同样是条件于时间配准结果和预设统计/工程阈值的判定。


\section{自动生成图表}
\input{generated/figures.tex}

\section{结论}
四个问题形成“同步--标定--融合--调度”的连续建模链。与直接将参考包算法整体移植不同，本论文保留当前仓库已经完成验证的 Q1--Q4 正式模型，仅吸收参考包中成熟的 LaTeX、图表自动引用、机器审计和提交打包能力。这样可以保证论文数值与正式结果文件同源，并降低人工复制造成的符号、数值和提交格式错误。

\appendix
\section{复现入口}
正式复现和交付流程统一由仓库根目录脚本驱动：
\begin{lstlisting}
python scripts/audit_inputs.py
python scripts/audit_results.py
python scripts/generate_latex_assets.py
python scripts/build_paper.py
python scripts/package_submission.py
\end{lstlisting}

\end{document}
